\chapter{Leib und Bewegung}

\section{Räumen und Zeitigen}

Fink hat den Vorrang, um den in der Tradition gestritten wird, zurückgewiesen, und zwar nach beiden Seiten. Kein Vorrang des Raumes, wie ihn die aristotelische Fundierung nahelegt, in der die Zeit als Zahl der Bewegung auf die Größe zurückgeht. Kein Vorrang der Zeit, wie ihn die Seelenlehre und dann Heideggers Fundamentalontologie behaupten. Stattdessen ein Vorrang der Welt, die als Zeit-Raum waltet, und ihr Walten ist die Bewegung: \zitat{Gesetzt den Fall, Bewegung sei ursprünglich das Räumen und Zeitigen}, dann müsste alles Bewegte in beidem sein, ohne dass eines das andere trüge.\footnote{Fink, Frühgeschichte, Kapitel zur Weltbewegung.} Die Bewegung ist bei Fink nicht ein Vorgang in Raum und Zeit, sondern das, was Raum und Zeit erst auseinanderlegt: Räumen wäre dann das Einräumen einer Weite, Zeitigen das Zumessen einer Weile.

Das ist ein Gedanke gegen alles, was bisher über den Rang der Zeit gesagt wurde, und er soll nicht verkleinert werden. Aber Fink selbst hält den Vorrang der Welt nicht ohne einen leisen Restvorrang der Zeit durch. Bei Platon liest er die Chora als das Raumhafte, die dunkle Materie, die den Ideen Platz gewährt, und den Nous als die Lichtmacht, die er in die Nähe der Zeit rückt; die Idee ist das Tun, die Chora das Leiden.\footnote{Fink, Frühgeschichte, Kapitel zu Platon.} Und die Bewegung, die er als Räumen und Zeitigen ansetzt, beschreibt er selbst so: \zitat{Bewegung ist die seltsame Weise, wo ein Ding im selben Moment an zwei Orten ist, d.\,h. im Übergang zwischen zwei Orten sich befindet.} Der Moment ist eines, die Orte sind zwei. Das Bewegte hält sich in einem Jetzt an zwei Orten auf, nicht an einem Ort zu zwei Zeiten. Das Räumen geschieht im Zeitigen, nicht das Zeitigen im Räumen. Finks Gleichrangigkeit hat also ihre eigene Richtung, und sie ist dieselbe wie die der anderen.

\section{Das Bewegte als Einheit}

Was Fink Welt nennt, nennt Hegel Materie und Bewegung. In der Nachschrift wird die Mechanik in zwei Teile geteilt: Raum und Zeit, und ihre Beziehung; die Beziehung hat zwei Gestalten, die Materie als relative Einheit und die Bewegung als absolute Einheit von Raum und Zeit.\footnote{Hegel, Naturphilosophie, Nachschrift, nach Bonsiepen.} Raum und Zeit können die Dialektik von Diskretheit und Kontinuität nicht an ihnen selbst verwirklichen; sie brauchen das, was in ihnen ist. Die Einheit von Zeit und Raum ist also nicht ein Drittes über ihnen, kein Behälter, in dem beide lägen, sondern das Bewegte, in dem beide zugleich sind.

Reichenbach sagt es mit der Genidentität: Längs der Zeit liegen die Zustände desselben Dinges, längs des Raumes verschiedene Dinge. Das Ding ist das, was die Faser bildet, und die Faser ist das, was die Zeit von der Zeit her und den Raum vom Nebeneinander her bestimmt. Ohne das Ding gäbe es weder Faser noch Nebeneinander, sondern nur einen Parameterraum, in dem nichts wirkt. Und die Parallelisierung von Raum und Zeit, die auf der Ebene der Mannigfaltigkeit so vollkommen erscheint, ist in Wahrheit die Abstraktion von dem, was in der Mannigfaltigkeit bewegt ist.

\section{Der Leib}

Unter den Bewegten ist eines ausgezeichnet, der Leib, und zwar nicht, weil er ein besonderer Körper wäre, sondern weil er das Bewegte ist, das sich selbst bewegt und das dabei jetzt und hier sagt. Koch hat gezeigt, dass er dem Raum die Richtungen gibt, Oben und Unten im Fallen, Vorn und Hinten im Gehen, Rechts und Links im Greifen. Baumann hat gezeigt, dass er die zweite Zeit gibt, die Zeit an Tag und Nacht, Wachen und Schlaf. Brentano hat gezeigt, dass er die dreidimensionale Grenze eines zeitlichen Kontinuums ist, das sich in ihm erneuert. Und Koch hat gezeigt, dass er das Partikulare ist, das sich a priori selbst individuiert und lokalisiert und damit die Spezifizierbarkeit aller anderen garantiert. Der Leib ist die Stelle, an der die Zeit im Raum ist: das Hier, das kein Dort ist, das Jetzt, das keine Konvention ist.

Das Verhältnis, das Koch als Wechselverhältnis von Subjektivität und raumzeitlicher Einzelnheit beschreibt, keine Seite ohne die andere, hat hier seinen Ort. Die Subjektivität braucht die raumzeitliche Einzelnheit, weil sie sonst kein Hier und kein Jetzt hätte, an dem sie sich individuiert. Die raumzeitliche Einzelnheit braucht die Subjektivität, weil sie sonst unspezifizierbar wäre, ein Universum von Doppelgängern. Und dieses Wechselverhältnis ist nicht symmetrisch, so wenig wie das von Zeit und Raum: Die Subjektivität ist der Ursprung des Werdens, die Einzelnheit sein Ergebnis.

\section{Die Weile im Fluss}

Guzzoni hat der Bewegung eine Gestalt gegenübergestellt, die die Tradition übersehen hat: das Bleiben. Die aristotelische Definition der Zeit als Zahl der Veränderung nach dem Früher und Später ist einseitig, denn nicht nur die Veränderung, auch die Ruhe hat Zeitcharakter. Die Weile ruht gleichsam im Fluss. Sie verläuft nicht linear, sie hat Anfang und Ende als eines zusammengesehen, sie ist offene Geschlossenheit. Und die Ruhe ist nicht das Fehlen von Bewegung, sondern ein \zitat{Raum für Bewegung}.\footnote{Guzzoni, Weile und Weite, Kapitel zur Weile und zur Ruhe.}

\begin{figure}[htbp]
\centering
\begin{tikzpicture}[scale=1,>=Stealth]
% Fluss der Zeit
\draw[->,very thick,decorate,decoration={snake,amplitude=1.2pt,segment length=9pt}] (-0.5,0) -- (10.5,0);
\node[below] at (10.2,-0.15) {\small Fluss der Zeit};
% Jetztpunkte
\foreach \x in {0.5,1.3,2.1,2.9} {\fill (\x,0) circle (0.05);}
\node[below] at (1.7,-0.2) {\scriptsize Jetzt, Jetzt, Jetzt \ldots};
% Weile als Fläche
\fill[gray!25] (4.2,-1.1) rectangle (7.6,1.1);
\draw[thick] (4.2,-1.1) rectangle (7.6,1.1);
\node at (5.9,0.55) {\small Weile};
\node at (5.9,-0.55) {\scriptsize keine halbe, keine doppelte};
% Ort als Augenblick des Raumes
\fill (5.9,0) circle (0.09);
\node[above right] at (5.95,0.05) {\scriptsize Ort};
\draw[<->,thin] (4.3,-1.35) -- (7.5,-1.35);
\node[below] at (5.9,-1.4) {\scriptsize Gegend};
% weiter Fluss
\foreach \x in {8.4,9.2} {\fill (\x,0) circle (0.05);}
\node[above] at (5.9,1.2) {\small \zitat{Diesem Fluß der Zeit gegenüber erscheint die Weile eher als eine Fläche oder ein Raum}};
\end{tikzpicture}
\caption{Die Weile im Fluss. Gegen die Folge der Jetztpunkte hat die Weile Fläche; in ihr ist der Ort der Augenblick des Raumes und die Gegend seine Weile (nach Guzzoni). Das Bleiben ist nicht Abwesenheit von Bewegung, sondern das, worin sie Raum hat.}
\label{abb:weile}
\end{figure}

Wenn Fink recht hat, dass die Bewegung Räumen und Zeitigen ist, dann hat Guzzoni recht, dass das Bleiben Einräumen und Weilen ist. Beide sind das Bewegte, in dem Zeit und Raum zugleich sind; das eine im Übergang, das andere im Aufenthalt. Und in beiden zeigt sich dieselbe Richtung. Im Übergang ist das Ding in einem Moment an zwei Orten: Die Zeit trägt, der Raum wird durchlaufen. Im Aufenthalt weilt die Zeit und gewinnt Fläche: Die Zeit wird Raum, der Raum nimmt sie auf. Es ist immer die Zeit, die etwas tut, und der Raum, dem es widerfährt. Finks Formel für Platon, Idee als Tun und Chora als Leiden, sagt es für den Ursprung des Gedankens. Aber Chora bedeutet auch das, ohne das nichts geschehen könnte. Das Leiden ist kein Mangel.
